%% kapa: exam=1 round=29 year=2018 period=1 form=A subject=civil_law subject_ko=민법 seq=21 number=21
%% source: Q-Net 감정평가사 기출문제 (cst00309) artlSeq=5202808 file=extracted/1008/5202808/2018년 제29회 감정평가사 1차시험 기출문제/2018년 제29회 감정평가사 1차시험 기출문제/1교시 A형.pdf
%% conversion: kordoc 4.24.51 -> tools/gen1.py (auto); review=proofread-sonnet; flags=none
\begin{question}{21}
지역권에 관한 설명으로 옳지 않은 것은?
\begin{choices}
  \item 지역권은 요역지와 분리하여 다른 권리의 목적으로 하지 못한다.
  \item 토지공유자의 1인은 지분에 관하여 그 토지를 위한 지역권 또는 그 토지가 부담한 지역권을 소멸하게 할 수 있다.
  \item 지역권자는 일정한 목적을 위하여 타인의 토지를 자기토지의 편익에 이용할 권리가 있다.
  \item 점유로 인한 지역권취득기간의 중단은 지역권을 행사하는 모든 공유자에 대한 사유가 아니면 그 효력이 없다.
  \item 계약에 의하여 승역지소유자가 자기의 비용으로 지역권의 행사를 위하여 공작물의 수선의무를 부담하기로 하고 이를 등기한 경우, 승역지소유자의 특별승계인도 그 의무를 부담한다.
\end{choices}
\end{question}
\answer{2}
